\documentclass[12pt]{article}
\usepackage[utf8]{inputenc}
\usepackage[T1]{fontenc}
\usepackage{hyperref}
\title{Yana Language Resources for K-12 Curriculum Development}
\author{Compiled by Claude Sonnet 5.5 (claude-sonnet-5-5) for Andr\'as Kornai}
\date{}
\begin{document}
\maketitle

\section{Introduction}
Yana (ISO 639-3 ynn) is an extinct language of north-central California. Sources record that it died out with Ishi in 1916, and no later fluent speaker is recorded for the Yahi dialect. Edward Sapir documented the dialects: Yahi from Ishi, Northern Yana through Betty Brown, and Central Yana through Sam Batwi. The documentation is substantial, and any curriculum would be heritage and revitalization work rather than instruction in a living language. Sources differ on the naming of the dialects.

\section{Primary Educational Resources}
No school materials or curricula for Yana were found in the search results.

\section{Linguistic and Reference Materials}
\begin{itemize}
\item \textit{Yana Dictionary} (University of California Press, 1960), by Sapir and Swadesh and edited by Mary Haas. A description from the UBC RelLex knowledgebase: \url{https://wiki.ubc.ca/Documentation:RelLex/Yana_Dictionary}
\item Wikipedia on the Yana language: \url{https://en.wikipedia.org/wiki/Yana_language}
\item Wikipedia on the Yahi dialect: \url{https://en.wikipedia.org/wiki/Yahi_dialect}
\item Wikipedia on Yana traditional narratives: \url{https://en.wikipedia.org/wiki/Yana_traditional_narratives}
\end{itemize}

\section{Religious and Community Resources}
A Yana-related page in the Sacred Texts collection appeared in the results (a mirror site): \url{https://neonvagabond.xyz/mirrors/sacred-texts/nam/ca/yat/index.htm}. Its content was not examined. No living community organization was identified in the search results.

\section{Digital and Audio Resources}
The American Philosophical Society holds the ACLS Yana materials: 26 notebooks, over 4,600 cards and slips, and 212 pages, including 25 notebooks Sapir recorded from Ishi with vocabulary files. The finding aid is at \url{https://indigenousguide.amphilsoc.org/entry/39835}. No audio was confirmed in the search results.

\section{Curriculum Development Context}
The archival record is the only source for the language. A heritage curriculum could combine the dictionary and Sapir's narratives with cultural and historical content. Any use of Ishi's materials should involve descendant and tribal communities, such as the Yahi's neighbours and the Native Californian groups connected to the region. These communities were not researched in this search.
\end{document}